\section{Datenquellen und Datenverarbeitung}

\subsection{Personendaten (5123101)}
\subsubsection{Transformation der Personendaten}
Als Grundlage für die Personendaten wurde ein Datensatz von Kaggle verwendet.\cite{swain_synthetic_person_data}

Die Rohdaten wurden zunächst mithilfe eines Python-Skripts transformiert und vereinheitlicht, um sie an die spezifischen Anforderungen des Projekts anzupassen. Im Zuge dieser Verarbeitung wurden folgende Anpassungen vorgenommen:

Zunächst wurde für jeden Datensatz eine \texttt{UUID} generiert, welche die ursprünglich verwendeten numerischen Identifikatoren ersetzt. Dadurch wird eine eindeutige Identifikation gewährleistet und die Verknüpfung von Einträgen über verschiedene Tabellen hinweg vereinfacht.

Der \texttt{Gender}-Wert wurde von ausgeschriebenen Bezeichnungen in eine standardisierte Kurzform (m, f, d) überführt. Die Attribute \texttt{Vorname, Nachname, PLZ, City, Street} und \texttt{StreetNo} wurden nach Bereinigung von überflüssigen Leerzeichen unverändert aus dem Datensatz übernommen. Für das Attribut \texttt{Country} wurde einheitlich der Wert „USA“ gesetzt, da sämtliche Adressdaten aus den Vereinigten Staaten stammen.

Da im Ausgangsdatensatz kein konkretes Geburtsdatum, sondern lediglich das Alter enthalten war, wurde zunächst das Geburtsjahr berechnet und anschließend durch ein zufällig generiertes Datum innerhalb dieses Jahres ergänzt. Auf diese Weise konnte für jeden Eintrag ein konsistenter \texttt{Birthdate}-Wert erzeugt werden.

Die vorhandenen E-Mail-Adressen entsprachen nicht den Anforderungen und wurden daher neu generiert. Hierfür wurden Kombinationen aus Vor- und Nachnamen sowie gängigen E-Mail-Domains verwendet. Abschließend wurden für das Attribut \texttt{Phonenumber} gültige Vorwahlen herangezogen und durch zufällig generierte Nummern ergänzt.

Das Ergebnis dieser Verarbeitung ist eine konsistent strukturierte und den Anforderungen entsprechende CSV-Datei, die als Grundlage für alle weiteren Verarbeitungsschritte dient.

\subsubsection{Aufteilung der Personendaten}
Im nächsten Schritt wurden die zuvor transformierten Personendaten auf die vorgesehenen Entitäten \texttt{Patient}, \texttt{Doctor} und \texttt{Nurse} verteilt beziehungsweise entsprechende Referenzen für die jeweiligen Tabellen erzeugt. Auch hierfür kamen Python-Skripte zum Einsatz.

Zunächst wurden alle \texttt{UUIDs} aus der transformierten \texttt{persons.csv} extrahiert. Anschließend wurden 95\,\% der Einträge der Entität \texttt{Patient} zugeordnet, während die verbleibenden 5\,\% als \texttt{Employees} klassifiziert wurden. Für jeden Patienteneintrag wurde zusätzlich eine numerische ID generiert, die als Primärschlüssel dient. Die entsprechenden Zuordnungen wurden anschließend in einer separaten CSV-Datei gespeichert.

Für die \texttt{Employees} wurde ebenfalls eine neue eindeutige ID in Form einer \texttt{UUID} erzeugt. Darüber hinaus erhielt jeder Eintrag ein zufällig zugewiesenes Department. Zusätzlich wurde das Attribut \texttt{currently\_working} eingeführt, welches bei 85\,\% der Einträge auf \texttt{true} gesetzt wurde, um aktive Mitarbeiter zu kennzeichnen. Diese Daten wurden ebenfalls in eine eigene CSV-Datei überführt.

Damit waren die Patientendaten vollständig aufbereitet. Im letzten Schritt erfolgte die weitere Differenzierung der \texttt{Employees} in \texttt{Doctors} und \texttt{Nurses}. Hierbei wurden 30\,\% der Mitarbeitenden als \texttt{Doctors} klassifiziert. Für diese Einträge wurden zusätzlich Telefonnummern mit einheitlicher Vorwahl generiert sowie ein spezifischer \texttt{Type} zugewiesen. Dabei wurde sichergestellt, dass jedes Department mindestens einen \texttt{head\_of\_department} enthält.

Die verbleibenden Mitarbeitenden wurden der Entität \texttt{Nurse} zugeordnet. Sowohl die Daten der \texttt{Doctors} als auch der \texttt{Nurses} wurden abschließend jeweils in separaten CSV-Dateien gespeichert.